\VerdierPhysicalPageAnchor{0114}%
\begingroup
\small
\setlength{\abovedisplayskip}{4pt plus 2pt minus 2pt}%
\setlength{\belowdisplayskip}{4pt plus 2pt minus 2pt}%

\VerdierSemanticAnchor{ch2-1-2}
\medskip
\noindent\textbf{1.2. Premières propriétés des catégories triangulées.}

\VerdierSemanticAnchor{ch2-1-2-1}
\medskip
\noindent\textbf{Proposition 1.2.1.}\quad
Soit $\mathcal D$ une catégorie triangulée. Pour tout objet $X$ de
$\mathcal D$, les foncteurs $\operatorname{Hom}_{\mathcal D}(X,\mathord\cdot)$
et $\operatorname{Hom}_{\mathcal D}(\mathord\cdot,X)$ sont des foncteurs
cohomologiques à valeurs dans la catégorie des groupes abéliens.

\VerdierSemanticAnchor{ch2-1-2-1-hom-functors}
\noindent\emph{Démonstration.}\quad
Il suffit tout d'abord de montrer que le foncteur
$\operatorname{Hom}_{\mathcal D}(X,\mathord\cdot)$ est un foncteur
cohomologique. L'autre assertion s'obtiendra alors en passant à la catégorie
opposée $\mathcal D^\circ$ (1.1.7). Soit donc

\VerdierSemanticAnchor{ch2-1-2-1-source-triangle}
\[
  X'\xrightarrow{u'}Y'\xrightarrow{v'}Z'\xrightarrow{w'}X'[1]
\]
un triangle distingué. Il suffit de montrer que dans le diagramme :

\VerdierSemanticAnchor{ch2-1-2-1-1}
\begin{equation}
  \resizebox{0.98\linewidth}{!}{$
    \operatorname{Hom}_{\mathcal D}(X,X')
    \xrightarrow{\operatorname{Hom}_{\mathcal D}(X,u')}
    \operatorname{Hom}_{\mathcal D}(X,Y')
    \xrightarrow{\operatorname{Hom}_{\mathcal D}(X,v')}
    \operatorname{Hom}_{\mathcal D}(X,Z')
    \xrightarrow{\operatorname{Hom}_{\mathcal D}(X,w')}
    \operatorname{Hom}_{\mathcal D}(X,X'[1])
  $}
  \tag{1.2.1.1}
\end{equation}

\VerdierSemanticAnchor{ch2-1-2-1-exactness-reduction}
on a :
\[
  \operatorname{Hom}_{\mathcal D}(X,v')
  \operatorname{Hom}_{\mathcal D}(X,u')=0
\]
et que ce diagramme est exact en $\operatorname{Hom}_{\mathcal D}(X,Y')$.
En effet, l'axiome TRII permettra alors d'obtenir le même résultat en
$\operatorname{Hom}_{\mathcal D}(X,Z')$. Soit donc $f\colon X\longrightarrow X'$.
D'après TRI, le triangle

\VerdierSemanticAnchor{ch2-1-2-1-identity-triangle}
\[
  X\xrightarrow{\operatorname{id}_X}X\longrightarrow0\longrightarrow X[1]
\]
est distingué et d'après TRII et TRIII, il existe un morphisme
$g\colon X\longrightarrow Y'$ tel que le diagramme ci-après soit commutatif :

\VerdierSemanticAnchor{ch2-1-2-1-2}
\begin{equation}
  \begin{tikzcd}[column sep=3.2em,row sep=2.5em]
    X \arrow[r,"{\operatorname{id}_X}"] \arrow[d,"f"']
      & X \arrow[r] \arrow[d,"g"']
      & 0 \arrow[r] \arrow[d]
      & X[1] \arrow[d,"{f[1]}"] \\
    X' \arrow[r,"u'"']
      & Y' \arrow[r,"v'"']
      & Z' \arrow[r,"w'"']
      & X'[1]
  \end{tikzcd}
  \tag{1.2.1.2}
\end{equation}

\VerdierSemanticAnchor{ch2-1-2-1-zero-composite}
On en déduit que $u'f=g$ et que $v'u'f=v'g=0$, ce qui démontre que
$\operatorname{Hom}_{\mathcal D}(X,v')
\operatorname{Hom}_{\mathcal D}(X,u')=0$.

\VerdierSemanticAnchor{ch2-1-2-1-lifting}
Soit maintenant $g\colon X\longrightarrow Y'$ un morphisme tel que
$v'g=0$. D'après TRII et TRIII, il existe un morphisme
$f\colon X\longrightarrow X'$ tel que le diagramme (1.2.1.2) soit
commutatif. On en déduit que $g=u'f$ et que par conséquent la suite
(1.2.1.1) est exacte en $\operatorname{Hom}_{\mathcal D}(X,Y')$.

\VerdierSemanticAnchor{ch2-1-2-2}
\medskip
\noindent\textbf{Corollaire 1.2.2.}\quad
Dans un triangle distingué
$X'\xrightarrow{u'}Y'\xrightarrow{v'}Z'\xrightarrow{w'}X'[1]$, le composé
de deux morphismes consécutifs est nul.

\VerdierSemanticAnchor{ch2-1-2-2-proof}
\noindent\emph{Démonstration.}\quad
En vertu de TRII, il suffit de montrer que $v'u'=0$. Pour cela, il suffit de
prendre $X=X'$ dans la suite exacte (1.2.1.1) et de considérer les images
successives de
$\operatorname{id}_{X'}\in\operatorname{Hom}_{\mathcal D}(X',X')$.

\VerdierSemanticAnchor{ch2-1-2-3}
\medskip
\noindent\textbf{Corollaire 1.2.3.}\quad
Soient $(X,Y,Z,u,v,w)$ et $(X',Y',Z',u',v',w')$ deux triangles distingués
de $\mathcal D$ et $(f_1,f_2,f_3)$ un morphisme de triangles. Si deux des
$f_i$, $1\leq i\leq3$, sont des isomorphismes, le troisième l'est aussi.
\endgroup

